\chapter{Fermions, Dirac Matrices, and the Dirac Sea}

The Dirac sea is our destination, but the route matters. Before facing negative energy in the relativistic electron theory, we need two pieces of machinery. First comes the operator algebra that distinguishes fermions from bosons. Then comes a review of the one-dimensional Dirac equation, where every enlargement of the wavefunction can be seen without the geometry of three spatial dimensions. Once those pieces are in place, the three-dimensional equation will force its own matrix structure, and the opening discussion of filled and empty fermion states will return as the key to positrons.

We use \(c=\hbar=1\). Capital \(P\) denotes a momentum operator, lower-case \(p\) a momentum eigenvalue, and bold symbols their three-dimensional versions.

\section{Exchange Becomes Operator Algebra}

Consider two identical particles, one at \(x\) and one at \(y\). The labels \(x\) and \(y\) denote complete positions; they are not two Cartesian coordinates of a single particle. Acting on the vacuum with two creation operators gives
\begin{equation}
\lvert x,y\rangle
=
\psi^\dagger(x)\psi^\dagger(y)\lvert0\rangle.
\label{eq:l10-two-particle-state}
\end{equation}

For bosons, exchanging the two slots leaves the state unchanged:
\begin{equation}
\lvert x,y\rangle=\lvert y,x\rangle.
\end{equation}
The operator statement that guarantees this symmetry is
\begin{equation}
\psi^\dagger(x)\psi^\dagger(y)
=
\psi^\dagger(y)\psi^\dagger(x).
\label{eq:l10-boson-exchange}
\end{equation}
Two bosonic creation operators commute, so we may reverse their order without changing the state they create.

For fermions, exchange changes the sign:
\begin{equation}
\lvert x,y\rangle=-\lvert y,x\rangle.
\label{eq:l10-fermion-state-exchange}
\end{equation}
The overall sign of an isolated state vector is not observable, but relative signs between alternatives are observable, so this minus sign must be carried consistently. Equation~\eqref{eq:l10-fermion-state-exchange} requires
\begin{equation}
\psi^\dagger(x)\psi^\dagger(y)
=
-\,\psi^\dagger(y)\psi^\dagger(x).
\label{eq:l10-fermion-operator-exchange}
\end{equation}

\begin{figure}[htbp]
\centering
\includegraphics[width=0.94\linewidth]{lecture_10_figure_02.png}
\caption{The board relation implementing fermionic exchange. Interchanging the creation operators reverses the sign of the two-particle state.}
\lectureframe{00:05:50}
\end{figure}

For any two operators \(A\) and \(B\), define the anti-commutator by
\begin{equation}
\{A,B\}:=AB+BA.
\end{equation}
The exchange rule is therefore
\begin{equation}
\{\psi^\dagger(x),\psi^\dagger(y)\}=0.
\label{eq:l10-creation-anticommutator}
\end{equation}
The corresponding fermionic field algebra replaces the bosonic commutators by anti-commutators. One must keep two uses of the word distinct: field operators anti-commute because of Fermi statistics, while matrices will later anti-commute in order to cancel cross terms in a relativistic dispersion relation.

Now set \(y=x\) in Eq.~\eqref{eq:l10-creation-anticommutator}. Both products become the same:
\begin{align}
0
&=\{\psi^\dagger(x),\psi^\dagger(x)\}\\
&=2\,\psi^\dagger(x)\psi^\dagger(x).
\end{align}
Hence
\begin{equation}
\bigl(\psi^\dagger(x)\bigr)^2=0.
\label{eq:l10-exclusion}
\end{equation}
Trying to create two identical fermions in the same one-particle state gives the zero vector. That is the exclusion principle in operator form.

\begin{figure}[htbp]
\centering
\includegraphics[width=0.94\linewidth]{lecture_10_figure_03.png}
\caption{The equal-state specialization on the board. The repeated creation-operator product appears twice and therefore must vanish.}
\lectureframe{00:08:04}
\end{figure}

\begin{figure}[htbp]
\centering
\begin{adjustbox}{max width=\linewidth}
\begin{tikzpicture}[node distance=1.25cm,>=Latex]
\node[draw,rounded corners,align=center,inner sep=7pt] (a)
  {$\{\psi^\dagger(x),\psi^\dagger(x)\}=0$};
\node[draw,rounded corners,align=center,inner sep=7pt,right=of a] (b)
  {$2(\psi^\dagger(x))^2=0$};
\node[draw,rounded corners,align=center,inner sep=7pt,right=of b] (c)
  {$(\psi^\dagger(x))^2=0$};
\draw[->,thick] (a) -- (b);
\draw[->,thick] (b) -- (c);
\end{tikzpicture}
\end{adjustbox}
\caption{The exclusion argument reconstructed without the partially erased intermediate board work.}
\end{figure}

\section{Why Filled and Empty Are Symmetric for a Fermion}

There is a second difference between bosonic and fermionic operators, subtler than exclusion and essential later. For a single bosonic mode,
\begin{equation}
[a,a^\dagger]=1,
\qquad
n=0,1,2,\ldots.
\end{equation}
Interchanging \(a\) and \(a^\dagger\) changes the sign of the commutator. Even if the operators were merely called \(A\) and \(B\), the algebra would tell us which ordering raises the occupation number and which lowers it.

The bosonic occupation ladder also has an orientation. The annihilation operator eventually reaches \(\lvert0\rangle\) and gives zero, whereas the creation operator can continue indefinitely:
\begin{align}
a\lvert0\rangle&=0,\\
a^\dagger\lvert n\rangle&\propto\lvert n+1\rangle,
\qquad n=0,1,2,\ldots.
\end{align}
Turning this ladder upside down does not reproduce the same spectrum. There is a bottom but no top.

For one fermionic mode,
\begin{equation}
\{a,a^\dagger\}=1,
\qquad
a^2=(a^\dagger)^2=0,
\qquad
n=0,1.
\label{eq:l10-fermion-mode}
\end{equation}
The creation operator takes the empty state to the filled state and then stops:
\begin{equation}
a^\dagger\lvert0\rangle=\lvert1\rangle,
\qquad
a^\dagger\lvert1\rangle=0.
\end{equation}
The annihilation operator performs the reverse:
\begin{equation}
a\lvert1\rangle=\lvert0\rangle,
\qquad
a\lvert0\rangle=0.
\end{equation}
There are only two states, so the diagram is symmetric under interchanging filled with empty and \(a\) with \(a^\dagger\). The operators still have different physical meanings, but their algebra no longer singles out one as creation and the other as annihilation. This formal symmetry is the seed of the particle-hole reinterpretation.

\begin{classroomqa}
\lecturetimestamp{00:13:54}
\audiencequestion{If \((a^\dagger)^2=0\), is that a contradiction, or does it imply that \(a^\dagger=0\)?}
\lecturerresponse{Neither. A nonzero operator can have a vanishing square. Here \(a^\dagger\) takes \(\lvert0\rangle\) to \(\lvert1\rangle\), but there is no allowed state above \(\lvert1\rangle\), so applying it twice gives zero. The same reasoning gives \(a^2=0\).}
\end{classroomqa}

Keep this filled-empty symmetry in reserve. When a negative-energy electron is removed, the annihilation operator will be relabeled as a positron creation operator. That move is not an arbitrary change of notation; it uses exactly the two-state fermion algebra established here.

\section{The One-Dimensional Dirac Construction}

Return now to one particle in one spatial dimension. The simplest first-order relativistic wave equation is
\begin{equation}
i\partial_t\psi=-\,i\partial_x\psi.
\label{eq:l10-right-mover}
\end{equation}
With \(P=-i\partial_x\), this is
\begin{equation}
H=P.
\end{equation}
A momentum eigenstate has \(E=p\), and its group velocity is
\begin{equation}
v_g=\frac{dE}{dp}=1.
\end{equation}
Positive or negative momentum, and therefore positive or negative energy, still gives a wave moving rigidly to the right. The sign of \(p\) changes the phase pattern, not the propagation direction.

A massive particle must be able to move in either direction. Introduce a second component satisfying
\begin{equation}
H\psi_1=P\psi_1,
\qquad
H\psi_2=-P\psi_2.
\label{eq:l10-two-directions}
\end{equation}
The first component moves right; the second moves left. Put them into a column,
\begin{equation}
\Psi=
\begin{pmatrix}
\psi_1\\
\psi_2
\end{pmatrix},
\end{equation}
and define
\begin{equation}
\alpha=
\begin{pmatrix}
1&0\\
0&-1
\end{pmatrix}.
\end{equation}
Both equations become
\begin{equation}
H\Psi=\alpha P\Psi.
\label{eq:l10-massless-one-d}
\end{equation}
The cost of including both directions is a new two-valued internal degree of freedom.

Direction alone does not produce mass. Introduce a second matrix,
\begin{equation}
\beta=
\begin{pmatrix}
0&1\\
1&0
\end{pmatrix},
\end{equation}
and try
\begin{equation}
H=\alpha P+\beta m.
\label{eq:l10-massive-one-d}
\end{equation}
The matrices satisfy
\begin{equation}
\alpha^2=\beta^2=\mathbb{1},
\qquad
\{\alpha,\beta\}=0.
\label{eq:l10-one-d-algebra}
\end{equation}
Squaring the Hamiltonian gives
\begin{align}
H^2
&=(\alpha P+\beta m)^2\\
&=\alpha^2P^2+\beta^2m^2
  +(\alpha\beta+\beta\alpha)Pm\\
&=P^2+m^2.
\label{eq:l10-one-d-dispersion}
\end{align}
The anti-commutator in Eq.~\eqref{eq:l10-one-d-algebra} removes the unwanted term linear in both \(P\) and \(m\).

The off-diagonal matrix \(\beta\) couples the two directional components. It is useful to picture the mass term as continually mixing right-moving and left-moving amplitudes. The picture should not be literalized into a particle bouncing back and forth along a microscopic track, but it captures why a massive wavepacket can have an average speed below light and can possess a rest frame, whereas an isolated massless chiral component cannot.

\begin{classroomqa}
\lecturetimestamp{00:23:37}
\audiencequestion{Is there a deeper physical interpretation of introducing both directional components?}
\lecturerresponse{There is, but it emerges more clearly in three dimensions. At this stage they are simply the two components required to describe opposite propagation and to admit an off-diagonal mass coupling. Their later interpretation is handedness, or chirality.}
\end{classroomqa}

\section{Three Dimensions Force the Pauli Matrices}

In three spatial dimensions momentum is a vector. A Hamiltonian linear in momentum and rotationally scalar therefore suggests another vector, now a vector whose components are matrices:
\begin{equation}
H=\boldsymbol{\alpha}\cdot\mathbf P
=\alpha_xP_x+\alpha_yP_y+\alpha_zP_z.
\label{eq:l10-three-d-ansatz}
\end{equation}

\begin{figure}[htbp]
\centering
\includegraphics[width=0.94\linewidth]{lecture_10_figure_04.png}
\caption{The three-dimensional massless ansatz \(H=\boldsymbol{\alpha}\cdot\mathbf P\) at the point where one matrix becomes a vector of three matrices.}
\lectureframe{00:26:00}
\end{figure}

Begin again with the massless case. We require
\begin{equation}
H^2=\mathbf P^2=P_x^2+P_y^2+P_z^2.
\end{equation}
Expanding the square of Eq.~\eqref{eq:l10-three-d-ansatz},
\begin{align}
H^2
&=\alpha_x^2P_x^2+\alpha_y^2P_y^2\\
&\quad+\alpha_z^2P_z^2\\
&\quad+\{\alpha_x,\alpha_y\}P_xP_y\\
&\quad+\{\alpha_x,\alpha_z\}P_xP_z\\
&\quad+\{\alpha_y,\alpha_z\}P_yP_z.
\end{align}
The diagonal terms are correct if
\begin{equation}
\alpha_x^2=\alpha_y^2=\alpha_z^2=\mathbb{1},
\end{equation}
and all cross terms disappear if
\begin{equation}
\{\alpha_i,\alpha_j\}=0,
\qquad i\ne j.
\label{eq:l10-alpha-anticommutation}
\end{equation}

Three familiar \(2\times2\) matrices do exactly this:
\begin{equation}
\sigma_x=
\begin{pmatrix}
0&1\\
1&0
\end{pmatrix},
\quad
\sigma_y=
\begin{pmatrix}
0&-i\\
i&0
\end{pmatrix},
\quad
\sigma_z=
\begin{pmatrix}
1&0\\
0&-1
\end{pmatrix}.
\label{eq:l10-pauli-matrices}
\end{equation}
Thus one solution is
\begin{equation}
\alpha_i=\sigma_i,
\qquad
H=\boldsymbol{\sigma}\cdot\mathbf P.
\label{eq:l10-weyl-hamiltonian}
\end{equation}

\begin{figure}[htbp]
\centering
\includegraphics[width=0.94\linewidth]{lecture_10_figure_05.png}
\caption{The compact massless Hamiltonian \(H=\boldsymbol{\sigma}\cdot\mathbf P\), with \(\alpha_z=\sigma_z\) still visible above it.}
\lectureframe{00:28:16}
\end{figure}

For a momentum eigenstate,
\begin{equation}
\boldsymbol{\sigma}\cdot\mathbf p
=
|\mathbf p|\,\boldsymbol{\sigma}\cdot\widehat{\mathbf p}.
\end{equation}
The operator \(\boldsymbol{\sigma}\cdot\widehat{\mathbf p}\) has eigenvalues \(+1\) and \(-1\), so
\begin{equation}
E=\pm|\mathbf p|.
\end{equation}
In this two-component theory, positive energy has the Pauli degree of freedom aligned with momentum, while negative energy has it anti-aligned. The alignment is locked; it is not an independently adjustable spin direction for a fixed energy sign.

\begin{classroomqa}
\lecturetimestamp{00:30:10}
\audiencequestion{Does spin drop out of the Dirac equation, or was it inserted because the electron is already known to have a magnetic moment?}
\lecturerresponse{Here “drop out” means “emerge as a consequence,” not “disappear.” The starting demand was a Hamiltonian linear in momentum whose square is \(\mathbf P^2\). Ordinary numbers cannot satisfy the required algebra, but the Pauli matrices can. Their two-component structure therefore arrives with the linearization itself. This does not prove that every particle must have spin; it says that particles described by this Dirac-type equation carry the corresponding spinor degree of freedom.}
\end{classroomqa}

This is one of the central lessons of the construction. The matrices are not decorative labels added after the relativistic equation is found. The attempt to take a local square root of the relativistic dispersion relation produces them.

\section{Why Mass Requires Four Components}

Try to repeat the one-dimensional mass construction without enlarging the \(2\times2\) theory:
\begin{equation}
H=\boldsymbol{\sigma}\cdot\mathbf P+\beta m.
\label{eq:l10-failed-mass}
\end{equation}
Assume \(\beta^2=\mathbb{1}\). The square contains the desired terms \(\mathbf P^2+m^2\), but it also contains
\begin{equation}
m\sum_i\{\sigma_i,\beta\}P_i.
\end{equation}
To remove this term, one matrix \(\beta\) would have to satisfy
\begin{equation}
\{\sigma_x,\beta\}
=
\{\sigma_y,\beta\}
=
\{\sigma_z,\beta\}
=0.
\label{eq:l10-impossible-beta}
\end{equation}
No \(2\times2\) matrix does so. The identity commutes with all three Pauli matrices, and every traceless \(2\times2\) matrix is a linear combination of the Pauli matrices. Any such combination fails to anti-commute with the component parallel to it. The \(2\times2\) algebra has no fourth independent matrix of the required kind.

The remedy is to enlarge the spinor. Four mutually anti-commuting matrices whose squares are the identity first appear in a \(4\times4\) complex representation. This matrix size is not chosen because spacetime has four dimensions; it is the smallest matrix space capable of realizing the required algebra.

There are many equivalent representations. A unitary change of spinor basis changes the entries of all the matrices without changing any physical prediction. What matters is the abstract set of relations
\begin{equation}
\{\alpha_i,\alpha_j\}=2\delta_{ij}\mathbb{1},
\qquad
\{\alpha_i,\beta\}=0,
\qquad
\beta^2=\mathbb{1}.
\label{eq:l10-dirac-algebra}
\end{equation}

\begin{classroomqa}
\lecturetimestamp{00:37:38}
\audiencequestion{Can \(\beta\) be regarded as a time component \(\alpha_t\)?}
\lecturerresponse{Not directly. The three \(\alpha_i\) do not combine with \(\beta\) as the four components of one Lorentz vector. In covariant notation one instead defines \(\beta=\gamma^0\) and \(\alpha_i=\gamma^0\gamma^i\); the four \(\gamma^\mu\) matrices satisfy the spacetime Clifford algebra.}
\end{classroomqa}

A convenient block representation is
\begin{equation}
\alpha_i=
\begin{pmatrix}
\sigma_i&0\\
0&-\sigma_i
\end{pmatrix},
\qquad
\beta=
\begin{pmatrix}
0&I_2\\
I_2&0
\end{pmatrix}.
\label{eq:l10-block-representation}
\end{equation}
The Pauli matrices in the diagonal blocks guarantee that different \(\alpha_i\) anti-commute. The off-diagonal identity blocks in \(\beta\) guarantee that \(\beta\) anti-commutes with every \(\alpha_i\).

\begin{figure}[htbp]
\centering
\includegraphics[width=0.94\linewidth]{lecture_10_figure_06.png}
\caption{The board's \(4\times4\) block representation. The upper equations record the required anti-commutators; the full matrices are typeset in Eq.~\eqref{eq:l10-block-representation}.}
\lectureframe{00:41:10}
\end{figure}

The massive Hamiltonian is now
\begin{equation}
H=\boldsymbol{\alpha}\cdot\mathbf P+\beta m
=
\begin{pmatrix}
\boldsymbol{\sigma}\cdot\mathbf P&mI_2\\
mI_2&-\boldsymbol{\sigma}\cdot\mathbf P
\end{pmatrix}.
\label{eq:l10-massive-dirac}
\end{equation}
Using Eq.~\eqref{eq:l10-dirac-algebra},
\begin{equation}
H^2=\mathbf P^2+m^2.
\end{equation}
The four-component structure is therefore not an optional embellishment. It is the smallest structure that permits both three-dimensional momentum and a mass term while keeping the equation first order.

\section{Chirality and the Meaning of the Extra Components}

What physical distinction did the enlargement add? Start with the massless upper block,
\begin{equation}
H_R=\boldsymbol{\sigma}\cdot\mathbf P,
\end{equation}
and the massless lower block,
\begin{equation}
H_L=-\boldsymbol{\sigma}\cdot\mathbf P.
\end{equation}
For positive-energy states, the first has spin aligned with momentum and the second has spin anti-aligned. These are opposite helicities. In a massless theory helicity is invariant under proper Lorentz transformations, because no observer can overtake a lightlike particle and reverse its momentum while leaving its spin unchanged. The two blocks therefore describe opposite chiral sectors.

\begin{classroomqa}
\lecturetimestamp{00:41:28}
\audiencequestion{In one dimension an extra binary degree of freedom was needed to introduce mass. What is the corresponding physical distinction in three dimensions?}
\lecturerresponse{It is the handedness of the spinor. The two \(2\times2\) blocks carry opposite relations between spin and momentum. One is right-handed and the other left-handed. The \(4\times4\) theory contains both.}
\end{classroomqa}

The mass term in Eq.~\eqref{eq:l10-massive-dirac} is off diagonal. It sends an upper-block amplitude into the lower block and vice versa:
\begin{equation}
\beta m
\begin{pmatrix}
\Psi_R\\
\Psi_L
\end{pmatrix}
=
m
\begin{pmatrix}
\Psi_L\\
\Psi_R
\end{pmatrix}.
\end{equation}
Mass is therefore a coupling between opposite chiralities.

\begin{classroomqa}
\lecturetimestamp{00:43:58}
\audiencequestion{Can a massless particle change chirality?}
\lecturerresponse{Not in the free massless Hamiltonian displayed here. The two blocks are decoupled, so a state that begins right-handed or left-handed remains so. The off-diagonal mass term is precisely what mixes them when \(m\ne0\).}
\end{classroomqa}

This language is exact for the block representation, but the image of a particle mechanically flipping between two arrows should not be pushed too far. The physical statement is the matrix coupling in the wave equation.

\begin{classroomqa}
\lecturetimestamp{00:44:50}
\audiencequestion{Are these matrices built into the quaternions?}
\lecturerresponse{The Pauli matrices are closely related to quaternion units: the matrices \(-i\sigma_i\) square to \(-\mathbb{1}\) and multiply with the quaternionic pattern, up to orientation conventions. The full Dirac matrices require the larger Clifford algebra, and the quaternion connection is not needed for the argument that follows.}
\end{classroomqa}

\begin{classroomqa}
\lecturetimestamp{00:46:34}
\audiencequestion{If spin is angular momentum, is there a separate energy term associated with it here?}
\lecturerresponse{No independent free-particle term is being added. The spinor structure already enters the kinetic Hamiltonian through the matrices. Coupling spin to an external magnetic field would introduce additional interaction terms, but that is a different problem.}
\end{classroomqa}

\section{The Alpha Matrices Are Velocity, Not Spin}

Each \(\alpha_i\) has two \(+1\) eigenvalues and two \(-1\) eigenvalues. Different components do not commute, so they cannot all be measured sharply at once. To identify what one component represents, use the Heisenberg equation. For an operator \(L\) without explicit time dependence,
\begin{equation}
\dot L=i[H,L].
\label{eq:l10-heisenberg}
\end{equation}
The sign convention is fixed here by the Schr\"odinger evolution \(e^{-iHt}\).

Apply Eq.~\eqref{eq:l10-heisenberg} to the position operator \(x\). The matrices act on spinor indices and commute with position, while only \(P_x\) fails to commute with \(x\):
\begin{align}
\dot x
&=i[\boldsymbol{\alpha}\cdot\mathbf P+\beta m,x]\\
&=i\alpha_x[P_x,x]\\
&=\alpha_x,
\label{eq:l10-velocity}
\end{align}
because \([P_x,x]=-i\). Thus \(\alpha_x\) is the \(x\)-component of the velocity operator; similarly,
\begin{equation}
\dot{\mathbf x}=\boldsymbol{\alpha}.
\end{equation}

\begin{classroomqa}
\lecturetimestamp{00:51:50}
\audiencequestion{Does a velocity-component measurement really give only \(+1\) or \(-1\)?}
\lecturerresponse{Yes, for the instantaneous operator \(\alpha_i\) in units \(c=1\). A general state need not be an eigenstate, so its expectation value can lie anywhere between those values. The smoothly transported wavepacket has an average velocity below light when the particle is massive.}
\end{classroomqa}

The velocity is not conserved:
\begin{equation}
\dot\alpha_x=i[H,\alpha_x]\ne0.
\end{equation}
The terms containing \(\alpha_y\), \(\alpha_z\), and \(\beta\) fail to commute with \(\alpha_x\). A state containing both positive- and negative-energy components consequently displays a rapid oscillatory contribution to its position. This is \emph{zitterbewegung}, literally trembling motion.

After averaging over times long compared with the rapid oscillation, the expected relativistic velocity is
\begin{equation}
\overline{\mathbf v}=\frac{\mathbf p}{E},
\end{equation}
which reduces to \(\mathbf p/m\) for a slow electron.

\begin{classroomqa}
\lecturetimestamp{00:54:40}
\audiencequestion{Does a free electron trembling without an external force violate Newton's law?}
\lecturerresponse{No. Translational invariance conserves momentum, so \(\dot{\mathbf P}=0\). The Dirac velocity operator is not simply \(\mathbf P/m\), and its rapid oscillatory part can change while momentum remains fixed. The time-averaged motion obeys the expected free-particle relation.}
\end{classroomqa}

\begin{classroomqa}
\lecturetimestamp{00:55:22}
\audiencequestion{Is the velocity always \(c\)?}
\lecturerresponse{An instantaneous measurement of one \(\alpha_i\) has eigenvalues \(\pm c\), but a massive wavepacket is generally a superposition and has a subluminal expectation value. Its mean group velocity is \(\mathbf p/E\).}
\end{classroomqa}

\begin{classroomqa}
\lecturetimestamp{00:55:30}
\audiencequestion{Can the oscillation be pictured as the mass varying the velocity?}
\lecturerresponse{No. The mass is a fixed parameter. The oscillation comes from noncommuting velocity and Hamiltonian operators, or equivalently from interference between energy sectors, not from a fluctuating mass.}
\end{classroomqa}

\begin{classroomqa}
\lecturetimestamp{00:56:22}
\audiencequestion{Could an instantaneous measurement find the velocity pointing one way even while the average motion points another way?}
\lecturerresponse{Yes. The component operator has the two sharp eigenvalues, while repeated or time-averaged measurements recover the smoother drift.}
\end{classroomqa}

The word spin must also be used carefully. The \(\alpha_i\) transform as components of velocity, an ordinary vector. The genuine spin operator in this representation is
\begin{equation}
\mathbf S=\frac12\boldsymbol{\Sigma},
\qquad
\Sigma_i=
\begin{pmatrix}
\sigma_i&0\\
0&\sigma_i
\end{pmatrix}.
\label{eq:l10-spin-operator}
\end{equation}
It differs from \(\alpha_i\) by the sign of the lower block. The matrices are closely related, but velocity and spin are not the same observable.

\section{Negative Energy and the Lowest-Energy Vacuum}

The negative-energy problem is already present before any of this matrix machinery. Return to the one-dimensional right mover,
\begin{equation}
H=P.
\end{equation}
Since \(p\) ranges over the whole real line, \(E=p\) has no lower bound.

The vacuum cannot merely mean a state with no particles. It must be the state from which no further energy can be removed. If an empty state allows us to add a negative-energy particle, lower the energy, add another, and continue indefinitely, then that empty state is not a stable vacuum.

Fermi statistics supplies a historical way out. Place one electron in every negative-energy state. Each added electron lowers the total energy, but exclusion allows only one electron in each one-particle state. Once every negative-energy state is filled, no additional electron can fall into that sector. Dirac identified this filled configuration as the vacuum.

The same construction would fail for bosons. Arbitrarily many bosons can occupy one state, so a negative-energy bosonic mode could absorb particles without limit. The fermionic bound \(n=0,1\), introduced at the beginning of the lecture, is exactly what stops the descent.

\begin{figure}[htbp]
\centering
\includegraphics[width=0.94\linewidth]{lecture_10_figure_07.png}
\caption{The historical Dirac-sea picture on the board: a ladder of negative-energy levels filled up to the dividing line, with positive-energy states above.}
\lectureframe{01:02:45}
\end{figure}

\begin{figure}[htbp]
\centering
\begin{adjustbox}{max width=\linewidth}
\begin{tikzpicture}[x=1.05cm,y=0.72cm,>=Latex]
\draw[->] (0,-4.8) -- (0,4.8) node[above] {$E$};
\draw[thick] (-2.9,0) -- (3.2,0) node[right] {$E=0$};
\foreach \y in {-4,-3.3,...,-0.5} {
  \draw (-2.3,\y) -- (2.3,\y);
  \foreach \x in {-1.8,-1.2,...,1.8} {
    \fill[blue!65!black] (\x,\y) circle (1.6pt);
  }
}
\foreach \y in {0.8,1.6,...,4.0} {
  \draw (-2.3,\y) -- (2.3,\y);
}
\draw[red,very thick] (0.6,-2.6) circle (5pt);
\draw[red,->,very thick] (0.6,-2.35) .. controls (1.4,-0.4) and (1.2,1.6) .. (1.5,2.4);
\fill[red] (1.5,2.4) circle (2.5pt);
\node[right,red] at (0.85,-2.6) {hole};
\node[right,red] at (1.7,2.4) {$e^-$};
\node[left] at (-2.5,-2.2) {filled negative-energy states};
\end{tikzpicture}
\end{adjustbox}
\caption{A clean reconstruction of excitation above a filled sea: promoting one negative-energy electron leaves a positive-energy hole.}
\end{figure}

The sea picture is historical language for a correct operator rearrangement. Its literal infinite charge and energy should not be mistaken for directly observable quantities; the modern reformulation will remove the need to imagine material electrons filling space.

\section{A Hole Is a Positron}

With all negative-energy states occupied, two kinds of excitation are possible. We may add a positive-energy electron above the sea. Or we may remove one negative-energy electron and leave a hole.

A hole has the opposite charge of the missing electron. Removing charge \(-e\) changes the charge by \(+e\). It also has positive excitation energy. If the removed electron had \(E_-<0\), then the energy measured relative to the filled vacuum changes by
\begin{equation}
E_{\mathrm{hole}}=-E_->0.
\end{equation}
The hole is a positive-energy, positively charged excitation: the positron.

A photon can promote a sea electron into a positive-energy state. Relative to the vacuum, the final state contains an electron and a positron. In a time-ordered diagram it is conventional to draw electron-number flow upward for the electron and the positron line with the opposite arrow. The arrows track fermion-number or charge flow; they are not little classical velocity vectors.

\begin{classroomqa}
\lecturetimestamp{01:05:14}
\audiencequestion{Does this picture require equal numbers of electrons and positrons?}
\lecturerresponse{Not for an arbitrary state. Starting from the vacuum and creating a neutral pair gives one of each, but a state may already contain excess electrons or excess positrons. Charge conservation constrains reactions; it does not require every allowed state to have equal populations.}
\end{classroomqa}

Dirac did not initially know that the hole would be a particle identical in mass to the electron and opposite only in charge and other additive quantum numbers. The equation and the hole construction led to that prediction before the positron was incorporated into the familiar particle catalogue.

\section{One Fermion Field Contains Electrons and Positrons}

The field operator brings the historical picture into modern algebra. Suppress spinor indices, normalization factors, and time dependence for the moment. A one-dimensional field may be expanded schematically as
\begin{equation}
\psi(x)
\sim
\int_{-\infty}^{\infty}dp\,a(p)e^{ipx}.
\label{eq:l10-full-field-integral}
\end{equation}
For the right-moving dispersion \(E=p\), positive \(p\) labels positive energy and negative \(p\) labels negative energy. Split the integral:
\begin{equation}
\psi(x)
\sim
\int_0^\infty dp\,a(p)e^{ipx}
+
\int_{-\infty}^0 dp\,a(p)e^{ipx}.
\end{equation}
In the second term set \(p=-q\), with \(q>0\), and reinterpret annihilation of a filled negative-energy electron as creation of a positron:
\begin{equation}
b^\dagger(q):=a(-q).
\end{equation}
Then
\begin{equation}
\psi(x)
\sim
\int_0^\infty dp\,
\left[
a(p)e^{ipx}
+
b^\dagger(p)e^{-ipx}
\right].
\label{eq:l10-electron-positron-split}
\end{equation}
The exact signs of the plane-wave phases depend on the spacetime convention; the invariant content is that \(\psi\) contains an electron annihilator and a positron creator.

\begin{figure}[htbp]
\centering
\includegraphics[width=0.94\linewidth]{lecture_10_figure_08.png}
\caption{The field integral split into positive- and negative-energy sectors. The lower line relabels the negative-energy electron operator as a positron creation operator.}
\lectureframe{01:11:35}
\end{figure}

This is where the filled-empty symmetry of Eq.~\eqref{eq:l10-fermion-mode} pays off. For fermions the operator algebra permits the relabeling without changing its anti-commutation structure.

A schematic electromagnetic interaction has the form
\begin{equation}
H_{\mathrm{int}}
\sim
\int dx\,\psi^\dagger(x)\psi(x)A(x),
\label{eq:l10-qed-schematic}
\end{equation}
where \(A\) denotes the photon field. Because each fermion field contains electron and positron pieces, a single local operator contains several channels:
\begin{align}
e^-&\longrightarrow e^-+\gamma,\\
e^-+e^+&\longrightarrow\gamma,\\
\gamma&\longrightarrow e^-+e^+,
\end{align}
together with their inverse and crossed processes. Energy-momentum conservation decides which channel is kinematically possible in a particular environment. For example, an isolated real photon in vacuum cannot create a pair without another participant to absorb recoil, even though the local vertex is present in the interaction.

\begin{figure}[htbp]
\centering
\includegraphics[width=0.94\linewidth]{lecture_10_figure_09.png}
\caption{The board juxtaposes electron emission, pair annihilation, and the schematic interaction \(\psi^\dagger\psi A\). Different selections from the same field operators generate the different channels.}
\lectureframe{01:13:30}
\end{figure}

One compact product of fields therefore organizes many apparently different reactions. Feynman later supplied the systematic diagrammatic rules for evaluating the amplitudes, but much of the operator structure is already visible here.

\section{Modern Vacuum Language and the Fermi-Sea Analogy}

\begin{classroomqa}
\lecturetimestamp{01:14:20}
\audiencequestion{How do modern physicists describe this without treating the vacuum as a literal Dirac sea?}
\lecturerresponse{They perform the operator reinterpretation directly. Negative-energy electron annihilation is renamed positive-energy positron creation, and the vacuum is defined as the state annihilated by all electron and positron annihilation operators. The resulting formulation is symmetric between particles and antiparticles and needs no material sea.}
\end{classroomqa}

The historical sea remains useful because its logic has a close condensed-matter analogue. In a crystal, a fixed number of electrons fill the lowest available one-particle states up to the Fermi level. The filled ground state is not empty. Promoting one electron above that level leaves an unoccupied state below it. Relative to the filled background, the hole carries positive excitation energy and can behave as a positive charge carrier.

There are important differences. A solid-state hole lives in a medium and its effective mass and charge response depend on a band structure. A positron is an elementary antiparticle in the relativistic quantum field. Even so, the common particle-hole algebra makes the analogy powerful.

\begin{classroomqa}
\lecturetimestamp{01:17:14}
\audiencequestion{Does the Klein--Gordon theory for bosons suffer the same cascade into negative-energy states?}
\lecturerresponse{No. A properly quantized complex Klein--Gordon field is organized into positive-energy particles and antiparticles without filling a bosonic sea. Its second-order equation and canonical structure differ from the Dirac Hamiltonian construction, so the historical fermionic sea argument is not transferred unchanged.}
\end{classroomqa}

\begin{classroomqa}
\lecturetimestamp{01:17:30}
\audiencequestion{Where does spin one for bosons come from?}
\lecturerresponse{Not from the Pauli-matrix linearization used here. Spin-one fields have their own Lorentz transformation law and field equations. The present derivation explains the spinor structure of Dirac fermions, not all possible particle spins.}
\end{classroomqa}

\begin{classroomqa}
\lecturetimestamp{01:17:46}
\audiencequestion{Why was Fermi statistics essential to the negative-energy solution?}
\lecturerresponse{Because exclusion caps every state at one electron. Once all negative-energy states are occupied, no further electron can fall into them. Bosonic occupation has no such cap, so a literal negative-energy boson sea would still have no bottom.}
\end{classroomqa}

\begin{classroomqa}
\lecturetimestamp{01:20:02}
\audiencequestion{Is electron versus positron simply another binary degree of freedom, like charge?}
\lecturerresponse{As state counting, a massive Dirac field has two electron spin states and two positron spin states, so two binary labels are a useful mnemonic. But particle versus antiparticle is not an independently rotatable internal qubit of one electron: every electron has charge \(-e\), every positron \(+e\), and the distinction emerges from the positive- and negative-frequency sectors after quantization.}
\end{classroomqa}

The infinite sea raises one final question.

\begin{classroomqa}
\lecturetimestamp{01:20:54}
\audiencequestion{Would the infinite Dirac sea produce an infinite cosmological constant?}
\lecturerresponse{It exposes the vacuum-energy problem. Bosonic zero-point modes contribute \(+\tfrac12\hbar\omega\), while fermionic modes contribute with the opposite sign. Formal cancellations are possible when spectra match, but the electron and photon spectra do not match, and an enormous remainder is expected. Without gravity, a constant shift of the Hamiltonian can be subtracted because only energy differences matter. With gravity, absolute vacuum energy gravitates, so the subtraction becomes a physical problem rather than mere bookkeeping.}
\end{classroomqa}

The last return is to the meaning of mass.

\begin{classroomqa}
\lecturetimestamp{01:22:56}
\audiencequestion{Can we say more about how mass relates left- and right-handed components?}
\lecturerresponse{The exact statement is already in the Hamiltonian: the off-diagonal term \(m\beta\) maps one chiral block into the other. Any more mechanical picture of microscopic flipping is optional and inevitably incomplete.}
\end{classroomqa}
